% Folder 138, batch 5 (archivists' pages 81-100).
% Pass: Opus 5.5 (claude-opus-5-5), 2026-09-30 — first pass, unchecked against the pages by a human.
%
% Pages 87 and 91 are blank and do not appear. Page 96 is an otherwise blank
% sheet carrying three lines in his hand: they become the heading of the run
% that follows (pp. 97-100) and the page gets no \page{}.
% Notation fixed for the batch: his curly T (the Teichmüller-type groups of
% pp. 83-89) is set \mathcal{T}; his script S (symmetric groups) \mathfrak{S};
% his script C (category of cartes, the torsor C of pp. 92-95) \mathcal{C};
% his double-struck letters \mathbb{}.
\documentclass[11pt,a4paper]{article}
\input{../preamble/grothendieck}

\folder{138}
\batch{5}
\pages{81}{100}
\foldertitle{Cartes. Etude arithmétique (1976, 77 ?) : bon de commande (s.d.), notes manuscrites (s.d.).}
\dating{[vers 1976-1978]}
\watermark{Édition de démonstration}

\begin{document}

\section{Action de Galois sur les systèmes $\Sigma(X/I/\xi/k)$}

\page{81}
\note{La page commence au milieu d'une phrase : l'argument vient d'avant le lot.}
des $\Phi_{i} \in \mathrm{Homext}(T, \pi)$ canoniques, $i \in I$. Le
système
\[
  \Sigma(X/I/\xi/k) = (T, I, \pi, (\Phi_{i}))
\]
\note{la notation $\Sigma(X/I/\xi/k)$ est ajoutée au-dessus d'une première
écriture biffée et illisible, $\struck{\Sigma(\ill{})}$ ; le troisième
argument est lu $\xi$ ici, \uncertain{$x$} n'est pas exclu.}
satisfait aux conditions envisagées. \struck{\ill{}} Au lieu \struck{\ill{}}
de prendre $x$, il suffisait de prendre un $\xi$, foncteur fibre sur la
catégorie des rev\supplied{êtements} étales de $X \setminus I$ ($\pi$ étant
le groupe des automorphismes de $\xi$\,\dots).

Alors $\Sigma$ dépend fonctoriellement de $(X/I/\xi/k)$. Prenant $X/I$, $k$
fixes, $\xi$ variable, on voit que $\Sigma$ est déterminé
\uncertain{« mod automorphismes intérieurs »} i.e. \struck{\ill{}} définit un
torseur \struck{sous} \add{à droite} \uncertain{non pas} sous
$\widetilde{A}_{\nu}$ mais sous son quotient
$(\widetilde{A}_{\nu}/\pi_{\nu} =)\ \widetilde{\Gamma}_{\nu}$,
\add{soit $P(X/I/k)$} \struck{(\ill{} dernier torseur} dont le commutant est
le groupe des « automorphismes ext. » de \struck{$(T, I, \pi$}
$\Sigma(X/I/\xi/k)$\,\dots

Supposons que $X/I/k$ provienne par extension de la base
\add{$k_{0} \to k$} \struck{\ill{} $S_{0}$} \struck{d'une situation}

\page{82}
\struck{$X_{0}/S_{0}$} d'une situation $X_{0}/I_{0}/k_{0}$, $k_{0}$ où
$k_{0}$ est un sous-corps. [OPS dans \uncertain{la} pratique $k$ est la
clôture \uncertain{alg.} de $k_{0}$\,\dots]. Alors
$\Gamma = \mathrm{Gal}(\overline{k}/k_{0})$ opère sur $\Sigma^{\mathrm{ext}}$
par transport de structure, donc \uncertain{aussi} une \struck{\ill{}}
opération \uncertain{associée} sur $P(X/I/k)$,
\add{\uncertain{i.e. on a}} \struck{donc} un homom.
\[
  \Gamma \longrightarrow {}^{P(X/I/k)}\widetilde{\Gamma}_{\nu}
  \simeq \mathrm{Aut}(\Sigma^{\mathrm{ext}})
\]
qui par \uncertain{changement} inverse \add{d'extension} donne naissance à
l'extension
\[
  1 \to \pi_{1}(X \setminus I, x) \to \pi_{1}(X_{0} \setminus I_{0}, x_{0})
  \to \pi_{1}(k_{0}, \overline{k}) \to 1
\]
\[
  \dots\qquad \Gamma \xrightarrow{\ \chi\ } \mathrm{Aut}(T)
  \simeq \widehat{\mathbb{Z}}^{*}
\]
\uncertain{est l'homom. de Kummer-Tate}\,\dots

\subsection*{II}

Plus \emph{précisément} supposons $k_{0} = \mathbb{Q}$,
$X_{0} = \mathbb{P}^{1}_{\mathbb{Q}}$, $\overline{\mathbb{Q}}$ clôture alg.
de $\mathbb{Q}$ dans $\mathbb{C}$,
\[
  I = \mu_{\nu+1}(\overline{\mathbb{Q}}) = \mu_{\nu+1}(\mathbb{C}^{*})
  \subset \overline{\mathbb{Q}} = \mathbb{E}^{1}_{\overline{\mathbb{Q}}}(\overline{\mathbb{Q}})
  \subset \mathbb{P}^{1}_{\overline{\mathbb{Q}}}(\overline{\mathbb{Q}})
\]
\note{$\mathbb{E}^{1}$ : lecture du symbole incertaine ; il désigne
manifestement la droite affine.}
On trouve un isomorphisme canonique, via l'exponentielle et les
\uncertain{lacets},
\[
  \struck{I \simeq [0,1]},\qquad
  \Sigma(X/I/k)^{\mathrm{ext}} \simeq L_{\nu}^{\mathrm{ext}}
\]
donc
\[
  P(X, I/k) = \text{torseur trivial}
\]

\page{83}
d'où
\[
  \Gamma = \mathrm{Gal}(\overline{\mathbb{Q}}/\mathbb{Q}) \longrightarrow \widetilde{\mathcal{T}}_{\nu}
\]
L'homom. correspondant
\[
  \Gamma \to \mathfrak{S}_{\nu+1}
\]
se factorise via
\[
  \Gamma \to \mu_{\nu+1}(\mathbb{C}^{*}) \simeq \mathbb{Z}/\nu{+}1
  \xrightarrow{\ \mathrm{can}\ } \mathfrak{S}_{\nu+1}\ \dots
\]

L'intérêt de prendre $I$ comme dessus est d'avoir une opération
\add{$\mathbb{Z}/2$} de $\mu_{\mathbb{Q}}$ \struck{\uncertain{transitive}}
sur la situation
\marginal{$n = \nu + 1$}
\[
  (X_{0}/I_{0}/\mathbb{Q})\quad \text{donc de } \overset{\mathbb{Z}/2}{\dots}\ \mu_{n}(\overline{\mathbb{Q}})\ (= D_{n})
\]
sur $X/I/\mathbb{Q}$ -- en fait, le produit semi-direct
$(\Gamma \times \struck{\ill{}}\ \mu_{n}(\overline{\mathbb{Q}}))
\rtimes \overline{\mathbb{Q}}^{*}$ opère sur $X/I/\overline{\mathbb{Q}}$ donc
\struck{\uncertain{j'aurais dû}}
\[
  (\Gamma \times \mathbb{Z}/2) \cdot \mu_{n}(\overline{\mathbb{Q}})
  \longrightarrow \widetilde{\mathcal{T}}_{\nu}
\]
\note{la ligne précédente, telle qu'écrite, place $\overline{\mathbb{Q}}^{*}$
en indice sous $\mu_{n}(\overline{\mathbb{Q}})$ ; la lecture comme produit
semi-direct est incertaine.}
où l'homom.
\[
  \mu_{n}(\overline{\mathbb{Q}}) \longrightarrow \widetilde{\mathcal{T}}_{\nu}
\]

\page{84}
se déduit de l'homom. évident
\[
  \mu_{n}(\overline{\mathbb{Q}}) \xrightarrow{\ \rho_{\nu}\ } \widetilde{A}_{\nu}
\]
provenant de la permutation circulaire des gén. de $\widehat{L}_{\nu}$ --
l'opération de $\mathbb{Z}/2$ dans $\widetilde{\mathcal{T}}_{\nu}$ provenant
de même de l'opération
\[
  \mathbb{Z}/2 \longrightarrow \widetilde{A}_{\nu}
\]
\marginal{vérifier ?}
par
\[
  (\lambda_{0}, \lambda_{1}, \dots, \lambda_{\nu}) \xrightarrow{\ t_{\nu}\ }
  \bigl(\lambda_{0},\ g_{\nu}\lambda_{\nu}g_{\nu}^{-1},\
  g_{\nu-1}\lambda_{\nu-1}g_{\nu-1}^{-1}, \dots,\ g_{2}\lambda_{2}g_{2}^{-1},\ \lambda_{1}\bigr)
\]
\note{au-dessous, une première écriture de l'image, biffée, en partie
illisible : $\struck{(\ill{}, \lambda_{\nu}, \lambda_{\nu-1}, \dots)}$.}
où $g_{i} = \lambda_{0} \cdots \lambda_{i-1}$, qui renverse l'ordre
circulaire.

Cela nous dit donc que l'homom.
\[
  \Gamma \xrightarrow{\ i_{\nu}\ } \widetilde{\mathcal{T}}_{\nu}
  \overset{\varepsilon_{\nu}}{\dashrightarrow}
  \mathfrak{S}_{n} \supset (\mathbb{Z}/n)^{*}
\]
\note{sous $(\mathbb{Z}/n)^{*}$, une addition biffée, lue
\uncertain{$\struck{\mathrm{Aut}(\mathbb{Z}/n)}$}.}
satisfait, en plus de
\[
  \varepsilon_{\nu}\, i_{\nu} = \chi_{n} \bmod n
\]
\note{encadré sur la page.}
la relation
\[
  \begin{cases}
    i_{\nu}(\Gamma) \text{ commute à } t_{\nu}\\[2pt]
    i_{\nu}(\gamma)\, \rho_{\nu}(i)\, i_{\nu}(\gamma)^{-1} = \rho_{\nu}(\chi_{n}(\gamma)\, i)
  \end{cases}
\]

\page{85}
Soit $\mathcal{T}_{\nu}^{!}$ le ss-groupe de $\widetilde{\mathcal{T}}_{\nu}$
formé des éléments qui satisfont aux relations précédentes. Donc on a
\[
  1 \to \mathcal{T}_{\nu}^{!\,0} \to \mathcal{T}_{\nu}^{!} \to
  \widehat{\mathbb{Z}}^{\times} \to 1
\]
avec
\[
  \mathcal{T}_{\nu}^{!\,0} \subset \mathcal{T}_{\nu}^{0},\qquad
  \mathcal{T}_{\nu}^{!\,0} =
  \Bigl\{\gamma \in \mathcal{T}_{\nu}^{0} \Bigm| \gamma \text{ commute aux } \rho_{\nu}(i)
  \text{ et à } t_{\nu}\Bigr\}
\]

\medskip
Si $\nu = 2$ il y a un autre choix intéressant de
$I_{0} \subset \mathbb{P}^{1}(\overline{\mathbb{Q}})$, savoir
$I = \{0, 1, \infty\}$ (plutôt que $I = \mu_{3}(\overline{\mathbb{Q}})$) qui
donne lieu à action de $\mathfrak{S}_{3}$ sur $(X_{0}/I_{0}/k)$ -- on trouve
action de $\Gamma \times \mathfrak{S}_{3}$ sur $\Gamma_{2}^{\mathrm{ext}}$,
i.e.
\[
  \Gamma \times \mathfrak{S}_{3} \longrightarrow \widetilde{\mathcal{T}}_{2}
\]
\note{sur la page, $\mathfrak{S}_{3}$ porte un accent au-dessus dans
« $\Gamma \times \mathfrak{S}_{3}$ », et $\Gamma_{2}^{\mathrm{ext}}$ est écrit
sur une première lettre surchargée ; lecture de ces deux points
incertaine.}

\page{86}
En fait on a ici, exceptionnellement
\[
  \widetilde{\mathcal{T}}_{2} \simeq \mathfrak{S}_{3} \cdot \mathcal{T}_{2}
\]
(produit semi-direct pour action de $\mathfrak{S}_{3}$ sur $\mathcal{T}_{2}$)
\[
  \Gamma \longrightarrow \struck{\ill{}}\ \widetilde{\mathcal{T}}_{2}
  \xrightarrow{\ \varepsilon\ } \mathfrak{S}_{3}\quad \text{trivial}
\]
i.e.\ $\Gamma \to \struck{\ill{}}\ \widetilde{\mathcal{T}}_{2}$, d'où
\[
  \boxed{\ \Gamma \longrightarrow \mathcal{T}_{2}^{\mathfrak{S}_{3}}\ }
\]

On aurait envie de définir un iso can.
$\mathcal{T}_{2}^{!} \simeq \mathcal{T}_{2}^{\mathfrak{S}_{3}}$
\marginal{?}
et des iso can.
\[
  \mathcal{T}_{2}^{!} \simeq \mathcal{T}_{\nu}^{!}
\]

Pour essayer une généralisation d'
\note{la page s'arrête sur cette phrase inachevée ; la page suivante (87) est
blanche.}

\section{\uncertain{Groupes profinis qui \ill{}}}

\page{88}
\note{titre souligné sur la page ; lecture incertaine au-delà de « Groupes
profinis ».}
\[
  \begin{cases}
    \pi \ \text{groupe profini}\\
    T\ \ \widehat{\mathbb{Z}}\text{-module libre de rang } 1
      \ \struck{\text{de } \widehat{\mathbb{Z}}\ (\text{isomorphisme non précisé})}\\
    \Phi_{i} \in \mathrm{Homext}(L, \pi)\ \struck{\ill{}}
      \qquad 0 \le i \le \nu \qquad (\nu \ge 2)
  \end{cases}
\]
\note{la deuxième ligne commence par $T$ (surchargé) ; la troisième écrit
$L$ là où l'on attend $T$ : transcrit tel quel.}

\textbf{Hyp}$_{\nu}$ :
\[
  \begin{cases}
    \exists\, e \in T^{*} \text{ base de } T\\
    (\varphi_{i})_{0 \le i \le \nu} \in \prod_{0 \le i \le \nu} \Phi_{i}
  \end{cases}
\]
tels que, posant $l_{i} = \varphi_{i}(e)$, l'homom. de
\[
  \widehat{L}_{\nu} \struck{\ill{}} \longrightarrow \pi,\qquad
  \lambda_{i} \longmapsto l_{i}
\]
\note{sous $\widehat{L}_{\nu}$ une glose partiellement biffée :
\struck{\uncertain{groupe profini \dots\ engendré par} $\lambda_{0}, \dots, \lambda_{\nu}$
\uncertain{avec} $\lambda_{0} \cdots \lambda_{\nu} = 1$}.}
soit un \emph{isomorphisme}.

En d'autres termes soit
\[
  \pi_{0} = \widehat{L}_{\nu} = \struck{\text{groupe profini}}\
  (\lambda_{0}, \dots, \lambda_{\nu} \mid \lambda_{0}\lambda_{1}\cdots\lambda_{\nu} = 1)^{\wedge}
\]
\[
  T_{0} = \widehat{\mathbb{Z}},\quad e_{0} = 1
\]
\[
  \Phi_{0i} \text{ classe de } \widehat{\mathbb{Z}} \longrightarrow \pi_{\nu},
  \qquad e_{0} \longmapsto \lambda_{i}
\]
on veut que $(\pi, T, (\Phi_{i}))$ soit isom.\ à
$(\pi_{0}, T_{0}, (\Phi_{0i}))$.
\[
  A_{\nu} = \mathrm{Aut}(\widehat{L}_{\nu}, T_{0}, (\Phi_{0i}))
  \xrightarrow{\ \chi_{\nu}\ } \widehat{\mathbb{Z}}^{*} = \mathrm{Aut}(T_{0})
\]
on prouve que c'est surjectif.
\[
  \widehat{L}_{\nu} \xrightarrow{\ i\ } \mathrm{Ker}\,\chi_{\nu} \subset A_{\nu}
  \quad\text{injectif},\qquad
  g \longmapsto (\mathrm{int}(g), \mathrm{id}_{T_{0}})
\]
$u\, i(g)\, u^{-1} = i(u(g))$ i.e.\ $\widehat{L}_{\nu}$ invariant dans
$A_{\nu}$.
\[
  A_{\nu}/\widehat{L}_{\nu} = \Gamma_{\nu}
\]
on veut élucider sa structure.

\page{89}
Considérons
\[
  L_{\nu} = \{\lambda_{0}, \dots, \lambda_{\nu} \mid \lambda_{0}\lambda_{1}\cdots\lambda_{\nu} = 1\}
  \quad \text{groupe discret}
\]
\[
  \Psi_{i} \in \mathrm{Homext}(\mathbb{Z}, L_{\nu})\ \text{classe de}\
  \psi_{i} : \mathbb{Z} \to L_{\nu},\quad \psi_{i}(1) = \lambda_{i}
\]
(donc $\Phi_{i} = \widehat{\Psi}_{i}$) \struck{\ill{}}

Considérons
\[
  B_{\nu} = \mathrm{Aut}(L_{\nu}, \mathbb{Z}, (\Psi_{i}))
\]
On a
\[
  B_{\nu} \xrightarrow{\ \chi_{\nu}^{0}\ } \mathbb{Z}^{*} = \{\pm 1\}
  \quad\text{surjectif}
\]
\[
  L_{\nu} \hookrightarrow \mathrm{Ker}(\chi_{\nu}^{0}),\qquad
  g \longmapsto (\mathrm{int}(g), \mathrm{id}_{\mathbb{Z}})
\]
Soit
\[
  \mathcal{T}_{\nu} = B_{\nu}/L_{\nu}
\]
(groupe de Teichmüller (discret) \uncertain{associé} pour la sphère à $\nu + 1$
trous)
\[
  1 \to \mathcal{T}_{\nu}^{0} \to \mathcal{T}_{\nu} \to \mathbb{Z}/2 \to 1
\]
On sait que $\mathcal{T}_{2}^{0} = 1$,
$\mathcal{T}_{\nu}^{0}$ = groupe des tresses sphériques colorées à $(\nu+1)$
brins, groupe non trivial si $\nu \ge 3$
$[\simeq \pi_{1}(\widetilde{D}_{\nu+1})\ ?]$.

On a
\[
\begin{tikzcd}
  \widehat{L}_{\nu} \arrow[r, hook] \arrow[d] & \widehat{B}_{\nu} \arrow[r, "\widehat{\chi}_{\nu}^{0}"] \arrow[d] & (\mathbb{Z}^{*})^{\wedge} = \pm 1 \arrow[d] \\
  \widehat{L}_{\nu} \arrow[r] & A_{\nu} \arrow[r, "\chi_{\nu}"] & \mathbb{Z}^{*}
\end{tikzcd}
\]
\note{la flèche verticale de gauche est tracée double ; les deux autres sont
des flèches d'inclusion. La dernière entrée de la ligne du bas est écrite
$\mathbb{Z}^{*}$ (et non $\widehat{\mathbb{Z}}^{*}$) : transcrit tel quel.}
\[
  \widehat{B}_{\nu}/\widehat{L}_{\nu} = \widehat{\mathcal{T}}_{\nu}
  \subset A_{\nu}/\widehat{L}_{\nu} = \Gamma_{\nu}
\]

\[
\begin{tikzcd}
  1 \arrow[r] & \widehat{\mathcal{T}}_{\nu}^{0} \arrow[r] \arrow[d, "\mathrm{can}"] & \widehat{\mathcal{T}}_{\nu} \arrow[r] \arrow[d] & \mathbb{Z}/2 \arrow[r] \arrow[d] & 1 \\
  1 \arrow[r] & \Gamma_{\nu}^{0} \arrow[r] & \Gamma_{\nu} \arrow[r] & \widehat{\mathbb{Z}}^{*} \arrow[r] & 1
\end{tikzcd}
\]

\emph{Questions} : $\widehat{\mathcal{T}}_{\nu}$ invariant dans
$\Gamma_{\nu}$ ? $\widehat{\mathcal{T}}_{\nu}^{0}$ dist.\ dans
$\Gamma_{\nu}^{0}$ ? Quand a-t-on \struck{\ill{}}
$\widehat{\mathcal{T}}_{\nu}^{0} = \Gamma_{\nu}^{0}$
(c'est faux pour $\nu = 2$).

\page{90}
\emph{Variante}. \uncertain{Donnons-nous} \ill{}
\[
  \varepsilon = (\varepsilon_{0}, \varepsilon_{1}, \dots, \varepsilon_{\nu})
  \in (\overline{\mathbb{N}^{*}})^{\nu+1}
\]
\note{sous $\overline{\mathbb{N}^{*}}$ une flèche renvoie à $[1, \infty]$.}
Soit
\[
  L_{\nu}^{\varepsilon} = \bigl\{\lambda_{0}, \dots, \lambda_{\nu} \bigm|
  \lambda_{0}\lambda_{1}\cdots\lambda_{\nu}
  = \lambda_{0}^{\varepsilon_{0}} = \lambda_{1}^{\varepsilon_{1}} = \dots
  = \lambda_{\nu}^{\varepsilon_{\nu}} = 1\bigr\}
\]
\note{la page s'arrête là ; la page 91 est blanche.}

\section{Le $\Gamma$-torseur $\mathcal{C}$}

\note{Pages 92--95 : autre encre (brune), autre feuillet.}

\page{92}
\[
  \left\{
  \begin{array}{l}
    \Gamma \text{ groupe},\ \mathcal{C}\ \Gamma\text{-torseur à droite}\\
    \pi,\ \Gamma\text{-groupe}\quad \struck{\ill{}}\\
    T_{i}\ (i \in I)\ (\pi, \Gamma)\text{-tors.}\\
    \pi_{i}\ \text{commutant}\\
    I_{i} \hookrightarrow \pi_{i}\ \text{sous-}\Gamma\text{-groupe}
  \end{array}
  \right.
  \qquad\Bigg|\qquad
  \begin{array}{l}
    \widetilde{\pi}\ \text{groupe}\\
    \widetilde{T}_{i}\ (i \in I)\ \widetilde{\pi}\text{-torseur}\\
    \widetilde{\pi}_{i}\ \text{son commutant}\\
    \widetilde{I}_{i} \subset \widetilde{\pi}_{i}
  \end{array}
\]
\[
  \mathcal{C} \xrightarrow{\ \rho\ } \mathrm{Isom}(\pi, \widetilde{\pi})
  \qquad \Gamma\text{-homomorphisme}
\]
\note{deux lignes barrées d'un zigzag au crayon :
$\struck{\ill{}\ \rho(\varphi) : \pi \to \widetilde{\pi}}$ et
$\struck{\ill{}\ T_{i} \wedge_{\pi} (\widetilde{\pi}, \varphi)}$.}
\[
  \pi \xrightarrow{\ \rho(\varphi)\ } \widetilde{\pi}
  \quad\text{donc}\quad \widetilde{\pi} \simeq \mathcal{C} \wedge_{\Gamma} \pi,
  \qquad
  T_{i} \xrightarrow{\ \rho_{i}(\varphi)\ } \widetilde{T}_{i}
\]
\[
  \begin{cases}
    \widetilde{\pi} = \mathcal{C} \wedge_{\Gamma} \pi\\
    (\widetilde{\pi}_{i}, \widetilde{T}_{i}, \widetilde{\pi}) = \mathcal{C} \wedge_{\Gamma} (\pi_{i}, T_{i}, \pi)\\
    \struck{\ill{}}\ \widetilde{I}_{i} \subset \widetilde{\pi}_{i}
      \qquad \widetilde{I}_{i} = \mathcal{C} \wedge_{\Gamma} I_{i}
  \end{cases}
\]
\note{à la fin de la troisième ligne : « donc $I_{i} = \dots$ », suivi d'une
formule surchargée et illisible, \struck{\ill{}}.}
\[
  \widetilde{I}_{i} \simeq \widehat{\mathbb{Z}}
  \quad\text{d'où}\quad
  \chi_{i} : \Gamma \to \widehat{\mathbb{Z}}^{*} \simeq \mathrm{Aut}(I_{i})
\]
\emph{indépendant de $i$}

\page{93}
\[
  \Gamma\ (= \mathrm{Aut}(\overline{\mathbb{Q}}))\ \text{groupe},\quad
  \mathcal{C}\ \Gamma\text{-tors.\ à droite}
  \quad (= \mathrm{Isom}(\overline{\mathbb{Q}}, \ill{}))
\]
\[
  \pi\ \bigl(= \pi_{1}(\mathbb{P}^{1}_{\overline{\mathbb{Q}}} - I_{\overline{\mathbb{Q}}})\bigr)
  \ \Gamma\text{-groupe}\quad
  (I \subset \mathbb{P}^{1}(\mathbb{Q})\ \text{fini})
\]
\[
  \forall i \in I\quad T_{i}\ \text{un}\ (\pi, \Gamma)\text{-torseur à droite}
  \quad (= \mathrm{Isom}(F_{\overline{\xi}}, F_{\overline{\xi}_{i}}))
\]
\[
  \forall\ r_{i} : \mathcal{C} \to T_{i}\ \text{satisfaisant}\quad
  r_{i}(\varphi\gamma) = \gamma^{-1} \cdot r_{i}(\varphi)
  \qquad \Big/\ r_{\cdot} : \mathcal{C} \to \textstyle\prod T_{i}
\]
\[
  \pi_{i} = \mathrm{Aut}_{\pi}(T_{i})\ \ \text{commutant de } \pi \text{ dans } T_{i}
  \ \ \text{ou}\ \ T_{i} \wedge_{\pi} (\pi\ \mathrm{int})
\]
\[
  I_{i} = T_{\cdot}(\overline{\mathbb{Q}})
  = \varprojlim_{n} \mu_{n}(\overline{\mathbb{Q}})
  \hookrightarrow \pi_{i}
\]
homom.\ de $\Gamma$-groupes avec $I_{i} \simeq \widehat{\mathbb{Z}}$
(pro-cyclique).

\medskip
\[
  \mathcal{C}_{2}(\pi\text{-ens}) \xrightarrow{\ f\ } (\mathrm{ens})
  \quad\text{\uncertain{ou} \ill{} de $\pi$}
\]
\[
  \begin{array}{l}
    \downarrow\ \overline{F}_{i}(E) = T_{i} \wedge_{\pi} E\\
    \pi_{i}\text{-}(\mathrm{Ens})\\
    \downarrow\ E_{i} \longmapsto E_{i}/I_{i} = \widehat{E}_{i}\\
    \mathrm{Ens}
  \end{array}
\]
\drawing{Une flèche courbe marquée $f_{i}$ va directement de
$\mathcal{C}_{2}(\pi\text{-ens})$ à $\mathrm{Ens}$ le long de la colonne ; un
trait fermé entoure le diagramme.}
\[
  \mathcal{C} \longrightarrow \mathrm{Isom}(\pi, \widetilde{\pi}).
\]

\page{94}
\[
\begin{tikzcd}
  \pi_{i} \arrow[r] & \pi \\
  I_{i} \arrow[u] \arrow[ur, "\alpha_{i}(\varphi)"'] &
\end{tikzcd}
\qquad I_{i} = I
\]
\[
  \alpha(\varphi\gamma)(\lambda) = \gamma_{\pi}^{-1}\,\alpha(\varphi)\bigl(\gamma_{I}(\lambda)\bigr),
  \qquad \gamma_{I}(\lambda) = \lambda^{\gamma}
\]
\[
  \gamma_{\pi}^{-1}\,\alpha_{i}(\varphi_{0})(\lambda_{i}) = l_{i}
\]
\[
  l_{1}, \dots, l_{n}\qquad
  \bigl(\gamma_{\pi}^{-1}(l_{1}^{\gamma}), \dots, \gamma_{\pi}^{-1}(l_{n}^{\gamma})\bigr)
\]
\[
  l_{1} l_{2} \cdots l_{n} = 1
\]
\[
  \gamma_{\pi}^{-1}(l_{1})
\]
\note{la page est coupée ici par un trait horizontal.}
\[
  \mathcal{C} \xrightarrow{\ \rho_{i}\ } \mathrm{Isom}(\widehat{\mathbb{Z}}, I_{i})
  \qquad \rho_{i}(\varphi)(1) = \lambda_{i}(\varphi)
\]
\[
  \pi \xrightarrow[\text{indépendant de } i]{\ \chi\ } \widehat{\mathbb{Z}}^{*}
  \qquad \lambda_{i}(\varphi\gamma) = \chi(\gamma)\,\lambda_{i}(\varphi)
\]
\[
  \rho_{i}(\varphi\gamma)(n) = \rho_{i}(\varphi)\bigl(\struck{n^{\chi(\gamma)}}\ \chi(\gamma) \cdot n\bigr)
  = \chi(\gamma)\,\rho_{i}(\varphi)(n)
\]
\[
  \lambda_{i}(\varphi) \in I_{i} \xrightarrow{\ \widetilde{r}_{i}(\varphi)\ } \pi
\]
\note{la flèche du haut du premier diagramme est lue de $\pi_{i}$ vers $\pi$ ;
à la fin de la dernière formule du bas, le $\rho_{i}(\varphi)(n)$ final est
surchargé.}

\page{95}
\[
  (\varphi \in \mathcal{C}) \longmapsto
  \pi \overset{\alpha_{i}(\varphi)}{\rightleftarrows} \pi_{i}
  \qquad \text{d'où}\quad I_{i} \hookrightarrow \pi_{i}
\]
\note{sous la double flèche une accolade : « bien \uncertain{compatible}
\ill{} \uncertain{avec} $\Gamma$ ».}
\[
  \alpha_{i}(\varphi\gamma) = \gamma_{\pi}^{-1}\,\alpha_{i}(\varphi)\,\gamma_{\pi_{i}}
\]
\[
  \mathcal{C} \xrightarrow{\ \beta\ } \mathrm{Isom}(\pi, \widetilde{\pi})
  \qquad \beta(\varphi\gamma) = \beta(\varphi) \circ \gamma_{\pi}
\]
\[
\begin{tikzcd}
  \pi \arrow[r, "\beta(\varphi)"] & \widetilde{\pi} \\
  I_{i} \arrow[u, "\alpha_{i}(\varphi)"] \arrow[ur, "\alpha'_{i}(\varphi)"'] &
\end{tikzcd}
\]
\[
  \begin{aligned}
    \alpha'_{i}(\varphi\gamma) &= \beta(\varphi\gamma)\,\alpha_{i}(\varphi\gamma)
      = \beta(\varphi)\,\gamma_{\pi}\,\gamma_{\pi}^{-1}\,\alpha_{i}(\varphi)\,\gamma_{\pi_{i}}\\
    &= \beta(\varphi)\,\alpha_{i}(\varphi)\,\gamma_{\pi_{i}} = \alpha'_{i}(\varphi)\,\gamma_{\pi_{i}}
  \end{aligned}
\]
donc
\[
  \widetilde{I}_{i} = \alpha'_{i}(\varphi)(I_{i})\quad \text{indépendant de } i
\]
\[
  \struck{\mathrm{Isom}(I_{i}, \widetilde{I}_{i}) =}\qquad
  \widetilde{I}_{i} = \mathcal{C} \wedge_{\pi} I_{i}
\]
\[
  \text{Si}\quad \struck{\pi' = \pi/\ill{} \dots\ \ill{}}\quad
  \pi' = \pi/U_{i},\quad U_{i} = \text{sous-groupe de } \pi \text{ opérant trivialement sur } I_{i}
\]
\note{« (indépendant de $i$) » est ajouté sous cette ligne.}
\[
  \widetilde{I}_{i} = \mathcal{C} \wedge_{\pi} I_{i}
\]
\[
  I_{i} \simeq T_{\cdot}(\overline{\mathbb{Q}})\qquad
  \pi' \simeq \mathrm{Aut}(I_{i}) \simeq \widehat{\mathbb{Z}}^{*}
\]
\[
  \mathcal{C} \simeq \varprojlim \mu_{n}(\overline{\mathbb{Q}})^{*}
\]
\note{l'indice « $\wedge_{\pi}$ » des lignes $\widetilde{I}_{i} = \mathcal{C}
\wedge_{\pi} I_{i}$ est écrit $\pi$ sur la page (et non $\Gamma$ comme en
p.~92) : transcrit tel quel.}

\section{Précartes et cartes toutes épinglées par précartes\,\dots}
\note{Titre repris des trois lignes de sa main, seules sur le feuillet
page~96 ; « épinglées » lu d'après « épi-glées » abrégé.}

\page{97}
Cartes \add{\uncertain{totalement}} épinglées etc.

\subsection*{I)}
Appelons \emph{pré-carte} un système $(S, A, F\,;\ p, q, r)$ où $S, A, F$
sont des ensembles, $p, q, r$ des applications
\[
  \struck{\ill{}\ S, A,}\qquad
  S \xrightarrow{\ p\ } \mathbb{N}^{*}\qquad
  A \xrightarrow{\ q\ } \mathbb{N}^{*}\qquad
  F \xrightarrow{\ r\ } \mathbb{N}^{*}
\]
et $q$ ne prenant que les valeurs 1 et 2.

Ils forment une catégorie pour \add{simplement} \struck{iso} \add{$\varphi$}
de $(S, A, F\,;\ p, q, r) \to{}$ $(S', A', F',$ $p', q', r')$
\struck{(\ill{} \dots\ \ill{})}\,\dots\
\emph{À toute carte combinatoire} $(\mathcal{C}, \dots)$ \ill{}
\struck{\uncertain{applications} pré-cartes}
\note{la ligne est entourée d'une accolade qui l'associe à la ligne
précédente ; lecture de l'ensemble incertaine.}
\[
  \struck{\varphi_{S} : S' \to S}
\]
\[
  \varphi_{s} : S' \to S,\qquad \varphi_{a} : A' \to A,\qquad \varphi_{f} : F' \to F
\]
satisfaisant
\[
  \left\{
  \begin{array}{l}
    p(s') = \nu(\varphi, s)\, p(s)\quad \struck{p(\varphi_{s}(s'))\ |\ p(\ill{})}\\
    q(a') = \nu(\varphi, a)\, q(a)\quad \struck{q(\varphi_{a}(a'))\ |\ q(\ill{})}\\
    r(f') = \nu(\varphi, f)\, r(f)\quad \struck{r(\varphi_{f}(f'))\ |\ r(\ill{})}
  \end{array}
  \right.
  \qquad
  \begin{array}{l}
    \forall s' \in S',\ s = \varphi_{s}(s')\\
    \forall a' \in A',\ a = \varphi_{a}(a')\\
    \forall f' \in F',\ f = \varphi_{f}(f')
  \end{array}
\]
\note{il a d'abord écrit des conditions de divisibilité, barrées, puis les
égalités au-dessus des lignes barrées.}
où $\nu(\varphi, s) = \mathrm{Card}(\varphi_{s}^{-1}(s))$,
$\nu(\varphi, a) = \dots$, $\nu(\varphi, f) = \dots$

On a un foncteur évident de la catégorie des cartes comb.\ \add{(Cat Comb)}
\uncertain{connexes}
\marginal{\uncertain{finies}}
\add{\uncertain{cartes ordinaires}} dans celle des précartes,
\add{(Précart)},
\[
  \mathcal{C} \longmapsto \mathrm{Fac}(\mathcal{C})
\]
\struck{Appelons épinglage d'} \add{\uncertain{C'est} \ill{}}
\struck{\ill{}}
le foncteur est \uncertain{conservatif}
\add{($\mathrm{Fac}(u)$ iso $\Rightarrow u$ iso)}
\marginal{D'où \emph{\uncertain{variété}} un foncteur \uncertain{fidèle}
\uncertain{sur} les iso}
Soit
\[
  \mathcal{C}_{0} = (S, A, F, p, q, r)\ \text{une précarte}
\]
\add{(\uncertain{connexe})}.
On appelle \emph{carte co-épinglée} \add{\uncertain{par $\mathcal{C}_{0}$}} la donnée
d'une carte $\mathcal{C}$ et d'un iso.
\[
  \struck{\mathrm{Fac}(\mathcal{C}) \to \mathcal{C}}\qquad
  \mathcal{C}_{0} \xrightarrow[\sim]{\ u\ } \mathrm{Fac}(\mathcal{C})
\]
\struck{C'est \ill{}} Les morphismes sont les morphismes de cartes
compatibles avec l'épinglage $u$, (\uncertain{ce sont}) \uncertain{bien}
\uncertain{des} isom.\ de cartes, et il a

\page{98}
trouve que la catégorie des cartes \add{orientées} co-épinglées est rigide --
donc équivalente à une catégorie discrète. \struck{D} Dans le cas où
$\mathcal{C}_{0}$ est fini, cette catégorie est elle-même finie. Elle est
\uncertain{vide} sauf si
\[
  \sum_{s} p(s) = \sum_{a} q(a) = \sum_{f} r(f)\ \ (= \delta),
\]
et si cette condition est vérifiée, \uncertain{alors} \struck{\ill{}} la
\add{\uncertain{cardinalité}} \struck{\ill{}} commune des \struck{\ill{}}
trois -- donc le \add{cardinal de l'} \struck{\uncertain{nombre des}}
classes d'isomorphie de co-cartes est égal à \uncertain{celui} des classes
d'iso.\ de $\mathfrak{S}_{2}^{0}$-ens \uncertain{de type}
« admissible » de ordre $\delta$, ayant les \ill{} \struck{\ill{}}
\uncertain{invariants} de \uncertain{ramification} que $\mathcal{C}_{0}$,
\emph{multiplié} par le cardinal du groupe des automorphismes de
$\mathcal{C}_{0}$, qui est égal à
\[
  \prod_{n \in \mathbb{N}^{*}} \bigl(a_{n}!\, b_{n}!\, c_{n}!\bigr)
\]
où $a_{n}, b_{n}, c_{n}$ sont les nb.\ de sommets (arêtes, faces) de degré
$n$.
\note{la parenthèse « (arêtes, faces) » est écrite telle ; il faut entendre
que $a_{n}$, $b_{n}$, $c_{n}$ comptent respectivement sommets, arêtes, faces.}

\marginal{//}
Définir ici $\mathcal{C}_{\mathcal{C}_{0}}$ et
$\mathcal{C}_{\mathcal{C}_{0}}^{\alpha}$, et l'opération de
$\Gamma \ni \gamma = \mathrm{Aut}(\mathcal{C}_{0})$ dessus. -- \textbf{NB}
$\mathcal{C}_{\mathcal{C}_{0}}^{\alpha}/\gamma$ est l'ens.\ des classes
d'iso.\ de cartes du type $(\mathcal{C}_{0}, \alpha)$\,\dots
\note{« $\Gamma \ni \gamma$ » : lecture incertaine ; on lit peut-être
$\Gamma_{\mathcal{C}_{0}}$.}

\page{99}
\subsection*{II)}
\add{\uncertain{Supposons} $\mathcal{C}_{0}$ fini. Soit un schéma de base $T$,}

\struck{(Appelons (revêtement} \struck{spécial} co-épinglé de
$\mathbb{P}^{1}_{T}$, un rev.\ fini de $\mathbb{P}^{1}_{T}$, étale sur
$U_{T} = \mathbb{P}^{1}_{T} - \{0, 1, \infty\}$, avec \struck{\ill{}}
\struck{\uncertain{isomorphisme}} \add{ramification} modérée fibre par fibre, et
\add{\uncertain{des} isomorphismes} \struck{des \ill{}}.
\note{ce paragraphe est barré de deux longues diagonales ; au-dessus de
« spécial » est ajouté $X \xrightarrow{f} \mathbb{P}^{1}_{T}$.}
\[
  \struck{f^{-1}(0_{T})_{\ill{}} \simeq}
\]
\[
  f^{-1}(0_{T})_{\mathrm{réd}} \simeq \struck{\ill{}}\ S_{T\,\mathrm{réd}}
\]
\[
  f^{-1}(1_{T})_{\mathrm{réd}} \simeq A_{T\,\mathrm{réd}}
\]
\[
  f^{-1}(\infty_{T})_{\mathrm{réd}} \simeq F_{T\,\mathrm{réd}}
\]
\note{après chaque $T$ un mot biffé, illisible, précède « réd ».}
\struck{\ill{} fibre par fibre sont \ill{}}
« compatibles avec les degrés de ramification »
\struck{\ill{}} \add{(\uncertain{partout}, \ill{})}
\marginal{\uncertain{fibre par fibre}, \uncertain{\ill{} ou partie} \ill{}}
On voit que cette catégorie est
rigide. \struck{Pour $\mathcal{C}_{0}$ fixé, soit}
\marginal{vérifier}
\struck{\ill{}} $F_{\mathcal{C}_{0}}(S)$ l'ens.\ des classes d'iso.\ de
cette catégorie. On voit que pour $S$ variable, on a un foncteur
contravariant en $S$, qui est représentable. On trouve que le schéma
$M_{\mathcal{C}_{0}}$ qui le représente est quasi-fini, et étale et séparé
sur $\mathbb{Z}$ ; il est de plus fini au-dessus de
$\mathbb{Z}[\frac{1}{\delta!}]$ -- i.e.\ aux voisinages des $p > \delta$
($\delta$ = complexité de $\mathcal{C}_{0}$).

\page{100}
On s'intéresse surtout au schéma $(M_{\mathcal{C}_{0}})_{\mathbb{Q}}$ -- qui
est donc fini étale sur $\mathbb{Q}$, non ramifié en les $p > \delta$.
\struck{\ill{} ce schéma \ill{}}

Sur le schéma $M_{\mathcal{C}_{0}}$, le groupe
$\gamma_{\mathcal{C}_{0}} = \mathrm{Aut}(\mathcal{C}_{0})$ opère de façon
évidente -- il opère sur $F_{\mathcal{C}_{0}}$. Le stabilisateur d'une
section de $M_{\mathcal{C}_{0}}$ sur un $T$, (i.e.\ d'une classe d'iso.\
de) $X \xrightarrow{f} \mathbb{P}^{1}_{S}$ \dots\ \struck{\ill{}}
\struck{\ill{}} co-\struck{\ill{}} épinglage, (si $S$ est connexe),
correspond \struck{aux} automorphismes du \struck{\ill{}}
rev.\ $X \xrightarrow{f} \mathbb{P}^{1}_{S}$ indépendamment de l'épinglage.
P.\ ex.\ si $\varphi$ est galoisien, il est can.\ iso.\ au groupe de Galois
de $f$.

P\uncertain{renons} \struck{Un} classe d'iso.\ $\alpha$ de
\struck{\uncertain{cartes}} \struck{(\uncertain{co-épinglées})} finis
\add{\uncertain{soit}} \struck{orientés (du type $\mathcal{C}_{0}$)},
$F_{\mathcal{C}_{0}}^{\alpha}$ le sous-foncteur \add{\uncertain{co-épinglé}}
\struck{\ill{}} $F_{\mathcal{C}_{0}}$ qui correspond aux
$X \xrightarrow{f} \mathbb{P}^{1}_{T}$ tels que les \add{\uncertain{ad}}
fibres géom., le groupe de Galois des
\note{la page s'arrête au milieu de la phrase : l'argument se poursuit
au-delà du lot.}

\end{document}
